\documentclass[11pt,a4paper]{article}
\usepackage{fontspec,amsmath,amssymb,mathtools}
\providecommand{\mathscr}[1]{\mathcal{#1}}
\DeclareFontFamily{U}{euf}{}
\DeclareFontShape{U}{euf}{m}{n}{<-> eufm10}{}
\usepackage[margin=22mm,headheight=14pt]{geometry}
\usepackage{graphicx,booktabs,longtable,array,calc,enumitem}
\usepackage{needspace,xurl,newunicodechar}
\usepackage[unicode,hidelinks]{hyperref}
\usepackage{bookmark,fancyhdr}
\defaultfontfeatures{Ligatures=TeX}
\usepackage{xeCJK}
\setmainfont{Times New Roman}
\setsansfont{Arial}
\setCJKmainfont[Path=../fonts/,
  BoldFont=NotoSansSC-SemiBold.ttf,
  ItalicFont=NotoSerifSC-Regular.ttf]{NotoSerifSC-Regular.ttf}
\setCJKsansfont[Path=../fonts/,
  BoldFont=NotoSansSC-SemiBold.ttf]{NotoSansSC-SemiBold.ttf}
\setCJKmonofont[Path=../fonts/]{NotoSansSC-SemiBold.ttf}
\xeCJKsetup{PunctStyle=quanjiao,CJKmath=true}
\renewcommand{\figurename}{图}
\XeTeXlinebreaklocale "zh"
\XeTeXlinebreakskip=0pt plus 1pt
\setmonofont{Consolas}[Scale=0.82]
\newfontfamily\extendedlatin{lmroman10-regular.otf}
\newunicodechar{ő}{{\extendedlatin\char"0151}}
\newunicodechar{ű}{{\extendedlatin\char"0171}}
\setlength{\emergencystretch}{4em}
\setlength{\parindent}{1.5em}
\setlength{\parskip}{2.5pt plus 1pt}
\linespread{1.10}
\allowdisplaybreaks[3]
\raggedbottom
\clubpenalty=10000
\widowpenalty=10000
\displaywidowpenalty=10000
\setcounter{secnumdepth}{-1}
\providecommand{\tightlist}{\setlength{\itemsep}{0pt}\setlength{\parskip}{0pt}}
\newcommand{\passthrough}[1]{#1}
\setlist{itemsep=2pt,parsep=0pt}
\makeatletter
\renewcommand\section{\@startsection{section}{1}{\z@}{-2.8ex plus -1ex minus -.2ex}{1.1ex plus .2ex}{\normalfont\large\bfseries}}
\renewcommand\subsection{\@startsection{subsection}{2}{\z@}{-2.1ex plus -.8ex minus -.2ex}{.8ex plus .2ex}{\normalfont\normalsize\bfseries}}
\renewcommand\subsubsection{\@startsection{subsubsection}{3}{\z@}{-1.8ex plus -.5ex minus -.2ex}{.6ex plus .2ex}{\normalfont\normalsize\bfseries}}
\makeatother
\pagestyle{fancy}
\fancyhf{}
\fancyhead[C]{\small 竞赛图中的 Hamilton 计数}
\fancyfoot[C]{\small\thepage}
\renewcommand{\headrulewidth}{0pt}
\hypersetup{pdftitle={本文结果及其与已有工作的关系},pdfauthor={简短说明}}
\begin{document}
\thispagestyle{plain}
\begin{center}
{\Large\bfseries 本文结果及其与已有工作的关系\par}
\vspace{9pt}
{\normalsize 简短说明\par}
\vspace{5pt}
{\small 2026 年 10 月 28 日\par}
\end{center}
\vspace{3pt}
\Needspace{5\baselineskip}
\section{1. 全局路径数的定量误差}\label{ux5168ux5c40ux8defux5f84ux6570ux7684ux5b9aux91cfux8befux5dee}

Divoux--Wang 已证明，竞赛图 Hamilton 路径极值的首项常数为 \(L=\cosh1/\cos1\)，奇阶圆环及其偶阶删点图渐近达到该值。本文保留这一常数与构造，证明

\[P(n)=\left(L+O(n^{-1})\right)\frac{n!}{2^{n-1}}.\]

证明直接从 Irving--Omar 所讨论的路径卷积出发，逐项控制永久式、删点、缩放及求和的误差。McCullagh 已有永久式的行列式近似和 \(O(n^{-1})\) 误差；Divoux--Wang 的证明在非负双随机情形也可细化到这一阶。这里的定量结论是经过全部计数步骤后，全图路径极值仍保留 \(O(n^{-1})\) 误差。

\Needspace{5\baselineskip}
\section{2. 原角函数的完整 Hessian}\label{ux539fux89d2ux51fdux6570ux7684ux5b8cux6574-hessian}

Wormald 的固定基图渐近式以角 Hessian 非退化为条件。本文从他的原隐式方程求导，对奇数 \(d\ge3\) 的正则基图得到

\[\mathbf H=E^T(I+C^2)(I-C^2)^{-1}E>0,
\qquad C=(A-A^T)/d.\]

这里 \(E=(e_1-e_d,\ldots,e_{d-1}-e_d)\)。压缩行列式中的因子 \(\det(E^TE)=d\) 给出

\[\sqrt{\frac d{\det\mathbf H}}
=\frac{\det(I+C)}{\det(I+iC)}=\rho_d(A-A^T).\]

正定性直接验证非退化条件，正有理行列式之比解释平方根的有理性；这就是对 Wormald Hessian 猜想的直接证明。 完整原角 Hessian、正定性、行列式/有理平方性质及圆环极限也见于 Bhattacharjee 的工作；这里给出详细推导和教学说明，不将这些重合结论列作相对于该文的新增结果。

\Needspace{5\baselineskip}
\section{3. 既有理论的应用与补充}\label{ux65e2ux6709ux7406ux8bbaux7684ux5e94ux7528ux4e0eux8865ux5145}

模型常数 \(\rho_d\) 本身可将期望加权邻接矩阵代入 Divoux--Wang 的一般计数定理得到。其谱上界与圆环闭式也见于他们的论文；正则内部块的 \(e^{1/d}\) 因子来自 Wormald。本文在这些结果上接入 Fang--Huang 的等号条件，以零行和确定圆环同构，再用非正则基图的指数衰减得到固定基图阶数下的期望唯一极值。严格单调性由已有显式频率推出，高阶展开由已知闭式计算。合稿将这些结论列作推论，保留具体证明并在相应位置注明来源。
\end{document}

